\thispagestyle{empty}

\vspace{0pt plus1fill} %число перед fill = кратность относительно некоторого расстояния fill, кусками которого заполнены пустые места
\begin{flushright}
  \large{\textit{На правах рукописи}}
%  \includegraphics[height=1.5cm]{my_signature} 
\end{flushright}

\vspace{0pt plus3fill} %число перед fill = кратность относительно некоторого расстояния fill, кусками которого заполнены пустые места
\begin{center}
\textbf {\large \MakeUppercase{Воробьев} Виталий Сергеевич}
\end{center}

\vspace{0pt plus3fill} %число перед fill = кратность относительно некоторого расстояния fill, кусками которого заполнены пустые места
\begin{center}
  \textbf{\Large \MakeUppercase{Модельно-независаимое}} \\
  \textbf{\Large \MakeUppercase{получение CP-нарушающих}} \\
  \textbf{\Large \MakeUppercase{параметров с использованием}} \\
  \textbf{\Large \MakeUppercase{когерентных состояний}} \\
  \textbf{\Large \MakeUppercase{нейтральных D-мезонов}}
%\textbf {\Large \thesisTitle}

\vspace{0pt plus3fill} %число перед fill = кратность относительно некоторого расстояния fill, кусками которого заполнены пустые места
{\large Специальность \thesisSpecialtyNumber\ "---\par <<\thesisSpecialtyTitle>>}

\vspace{0pt plus1.5fill} %число перед fill = кратность относительно некоторого расстояния fill, кусками которого заполнены пустые места
{\Large\MakeUppercase{Автореферат}}\par
\large{диссертации на соискание учёной степени\par \thesisDegree}
\end{center}

\vspace{0pt plus4fill} %число перед fill = кратность относительно некоторого расстояния fill, кусками которого заполнены пустые места
\begin{center}
{\large{\thesisCity\ "--- \thesisYear}}
\end{center}

\newpage
% оборотная сторона обложки
%\thispagestyle{empty}
\noindent Работа выполнена в \thesisInOrganization.

\par\vspace{2mm}%\bigskip
%\begin{table}[h] % считается не очень правильным использовать окружение table, не задавая caption
    \noindent%
    {\sfs \MakeUppercase{Научный руководитель}:}\\
%     
%     \vspace{-5 mm}
    \noindent
    \begin{tabular}{@{}llp{10.8cm}}
      \begin{tabular}[t]{@{}l@{}}
        \textbf{}\MakeUppercase{Бондарь}\\
        \textbf{}{Александр Евгеньевич}
      \end{tabular}
      & --- & {\sfs \supervisorRegalia, \par
      \thesisOrganizationShort, г.~Новосибирск.%, зам. директора
        %\par
        \vspace{3mm} 
        }
      \end{tabular}  
      
      \noindent
      {\sfs \MakeUppercase{Официальные оппоненты}:} \\% \vspace{2 mm}
%      \noindent\vspace{-4 mm}
      \noindent\begin{tabular}{@{}llp{10.8cm}}
        \begin{tabular}[t]{@{}l@{}}
         \textbf{}\MakeUppercase{Николаенко}\\
         \textbf{}{Владимир Иванович}
        \end{tabular} 
        & --- &  {\sfs \opponentOneRegalia,\par
                  \opponentOneJobPlace, %\par
                  \opponentOneJobPost%\par 
                  \vspace{2mm}
        } \\
        \begin{tabular}[t]{@{}l@{}}
         {\sfs\textbf{}\MakeUppercase{Ростовцев}}\\
         {\sfs\textbf{}{Андрей Африканович}}
        \end{tabular} 
        & --- &
        {\sfs     \opponentTwoRegalia,\par
                  \opponentTwoJobPlace, %\par
                  \opponentTwoJobPost
        }
%        \vspace{1mm}
        \\
        \begin{tabular}[t]{@{}l@{}}
         {\sfs \MakeUppercase{Ведущая}}\\
         {\sfs \MakeUppercase{организация}:}
        \end{tabular}
        & --- &
         {\sfs \leadingOrganizationTitle }
    \end{tabular}  
%\end{table}
\par\vspace{2mm}%\bigskip

%\noindent 
Защита состоится \defenseDate~на~заседании диссертационного совета \defenseCouncilNumber~на базе \defenseCouncilTitle~по адресу: \defenseCouncilAddress.

\vspace{1mm}
%\noindent 
С диссертацией можно ознакомиться в библиотеке \synopsisLibrary.

\vspace{1mm}
%\noindent
{Автореферат разослан \synopsisDate.}

%\vspace{1mm}
%\begin{table} [h] % считается не очень правильным использовать окружение table, не задавая caption
%\par\bigskip
    \noindent\hspace{-6mm}
    \begin{tabular}{p{13cm}cr}
        \begin{tabular}{p{10cm}}
            {\sfs Ученый секретарь}  \\
            {\sfs диссертационного совета}  \\
%            \sfs \defenseCouncilNumber, 
            {\sfs \defenseSecretaryRegalia}
        \end{tabular} 
    &
        \begin{tabular}{c}
%            \includegraphics[height=1cm]{secretary-signature} 
        \end{tabular} 
    &
        \begin{tabular}{r}
            \\
            \\
            {\sfs \defenseSecretaryFio}
        \end{tabular} 
    \end{tabular}
%\end{table}
\newpage
